\documentclass[11pt]{article}
\usepackage[margin=1in]{geometry}
\usepackage[T1]{fontenc}
\usepackage[utf8]{inputenc}
\usepackage{lmodern}
\usepackage{graphicx}
\usepackage{booktabs}
\usepackage{array}
\usepackage{amsmath}
\usepackage{hyperref}
\usepackage{parskip}
\usepackage{xcolor}

\hypersetup{
  colorlinks=true,
  linkcolor=blue!50!black,
  urlcolor=blue!50!black
}

\title{Citation Imputation Methods Note}
\author{}
\date{\today}

\begin{document}
\maketitle

\section*{Purpose}
This note isolates the imputation problem from the policy counterfactual. The
question here is narrower: \emph{given the papers with observed OpenAlex
citations, which imputation model best predicts downstream citations for the
2018--2020 RDD sample?}

\section*{1. Evaluation Setup}
The evaluation set is the common observed-paper subset where all methods are
comparable: 1{,}511 papers from the 2018--2020 RDD bandwidth sample with
observed citations, abstract embeddings, and full-text embeddings.

Observed citations are evaluated in both log and level space. The main ranking
metric is Spearman correlation on $\log(1+\text{citations})$; the main error
metric is mean absolute error in the same space. This reflects the actual use
case: we care more about preserving relative ordering than about exact
paper-level citation counts.

All neighbor-based methods are year-aware: a paper can only borrow information
from papers in the same or earlier year.

\section*{2. Methods}
\textbf{Abstract $K=20$ kNN.}
This is the original baseline. Each paper is represented by a single SPECTER2
embedding of \texttt{title [SEP] abstract}. Prediction is a similarity-weighted
average of the 20 nearest observed papers in cosine space.

\textbf{Full-text $K=20$ kNN.}
This keeps the same neighbor rule but upgrades the representation. For each
paper we remove the \texttt{References} section, split the full submission body
into evenly spaced chunks, embed each chunk with SPECTER2, and average the
chunk vectors into one paper embedding.

\textbf{Full-text topic-restricted $K=20$ kNN.}
This uses the same full-text representation but restricts neighbors to the same
coarse topic when enough same-topic neighbors exist. When the topic pool is too
small, the method falls back to the global pool.

\textbf{Full-text ridge.}
This uses the same full-text representation as the two full-text kNN methods,
but replaces the neighbor rule with a supervised ridge regression on the
embedding coordinates. This is still year-aware in evaluation: each paper is
predicted using papers from the same or earlier year only.

\section*{3. Main Benchmark Results}

\begin{center}
\small
\begin{tabular}{>{\raggedright}p{4.2cm}rrrr}
\toprule
Method & Pearson(log) & Spearman(log) & MAE(log) & MAE(level) \\
\midrule
Abstract $K=20$ kNN & 0.246 & 0.216 & 1.125 & 68.7 \\
Full-text $K=20$ kNN & 0.235 & 0.216 & 1.125 & 68.7 \\
Full-text topic $K=20$ kNN & 0.221 & 0.205 & 1.127 & 68.8 \\
Full-text ridge & 0.321 & 0.299 & 1.081 & 67.1 \\
\bottomrule
\end{tabular}
\end{center}

Three conclusions are clear.

First, moving from abstracts to full submission text does not help by itself.
The two full-text nearest-neighbor variants do not beat the original abstract
baseline.

Second, the improvement comes from switching from a neighbor rule to a
supervised model on the same full-text representation. Full-text ridge is the
only method that clearly improves both ranking and error.

Third, the gain is not small. Relative to the original abstract baseline, ridge
raises Spearman from 0.216 to 0.299 and lowers log-space MAE from 1.125 to
1.081.

\begin{figure}[p]
  \centering
  \includegraphics[width=0.90\textwidth]{figures/model_comparison.png}
  \caption{Observed-paper comparison of the four imputation methods. Full-text
  ridge is the clear winner.}
\end{figure}

\begin{figure}[p]
  \centering
  \includegraphics[width=0.98\textwidth]{figures/predicted_vs_actual.png}
  \caption{Predicted versus actual $\log(1+\text{citations})$ for each method.
  The ridge model still compresses the outcome, but less severely than the
  neighbor-based methods.}
\end{figure}

\section*{4. Neighbor Sensitivity: $K=5$ vs. $K=20$}
We also reran the neighbor-based methods with $K=5$ instead of $K=20$.

\begin{center}
\small
\begin{tabular}{>{\raggedright}p{4.0cm}rrrr}
\toprule
Method & Spearman($K=20$) & Spearman($K=5$) & MAE($K=20$) & MAE($K=5$) \\
\midrule
Abstract kNN & 0.216 & 0.180 & 1.125 & 1.216 \\
Full-text kNN & 0.216 & 0.179 & 1.125 & 1.216 \\
Full-text topic kNN & 0.205 & 0.190 & 1.127 & 1.202 \\
\bottomrule
\end{tabular}
\end{center}

Reducing the neighbor count makes every nearest-neighbor imputer worse. The
small local neighborhood is noisier, and the extra locality does not buy enough
precision to compensate. So if we keep a nearest-neighbor baseline for
comparison, $K=20$ is the better operating point than $K=5$.

\section*{5. Recommendation}
The recommended default imputer is now \textbf{full-text ridge}. It is the only
method that materially improves on the original baseline, and it does so with a
simple and transparent model. The full-text kNN variants remain useful as
diagnostic baselines, but they are no longer the main path for policy or
counterfactual work.

\end{document}
